\selectlanguage{english}
\begin{abstract}
First- and second-order models with time delay were used to
model the temperature behaviour of an electric furnace. The first-order
methods included Alfaro's two-point method and point-averaging-based models.
For the second-order methods, the approaches of Jahanmiri and Fallahi, Stark’s method,
and first- and second-order models generated in Matlab were applied.

The results indicate that the Matlab-generated models provided the most accurate
predictions of system behaviour, while the first-order methods exhibited lower error
compared to the second-order methods.
\end{abstract}

\begin{IEEEkeywords}
System identification; Alfaro's method; Stark's method; Jahanmiri and Fallahi method.
\end{IEEEkeywords}

\selectlanguage{spanish}
\begin{abstract}
Se emplearon métodos de primer y segundo orden con retardo
para modelar el comportamiento de la temperatura de un horno eléctrico.
Los métodos de primer orden incluyeron el método de dos puntos de Alfaro
y modelos basados en el promedio de puntos. Para los métodos de segundo orden,
se aplicaron los enfoques de Jahanmiri y Fallahi, el método de Stark,
y los modelos de primer y segundo orden generados en Matlab.

Los resultados indican que los modelos generados en Matlab proporcionaron
las predicciones más precisas del comportamiento del sistema,
mientras que los métodos de primer orden presentaron un error menor en comparación
con los métodos de segundo orden.
\end{abstract}

\begin{IEEEkeywords}
Identificación de sistemas; método de Alfaro; método de Stark; método de Jahanmiri y Fallahi.
\end{IEEEkeywords}
